% Template for ICIP-2024 paper; to be used with:
%         spconf.sty  - ICASSP/ICIP LaTeX style file, and
%         IEEEbib.bst - IEEE bibliography style file.
% --------------------------------------------------------------------------
\documentclass{article}
\usepackage{spconf,amsmath,graphicx, amssymb}
\usepackage{amsfonts} % For \mathbb

% Example definitions.
% --------------------
\def\x{{\mathbf x}}
\def\L{{\cal L}}

% Title.
% ------
\title{Pragmatic Learning: Intrinsic communicative clarity in multi-agent settings}
%
% Single address.
% ---------------
\name{Henri Langlois, Pablo Piantanida}
\address{CNRS, ILLS, Mila, CentraleSupelec}

\begin{document}
%\ninept
%
\maketitle
%
\begin{abstract}
]Reinforcement Learning (RL) agents often fail in cooperative settings because their actions, while individually optimal, are unpredictable to their partners. This challenge is magnified in human-agent teams, where an agent's opacity can erode trust and team effectiveness. We address this limitation by first establishing a formal equivalence between the principles of pragmatic inference, as formalized by the Rational Speech Act (RSA) framework, and Maximum Entropy Reinforcement Learning (Max-Ent RL). This theoretical bridge allows us to derive a principled intrinsic reward that incentivizes "legible" behavior—actions that clearly communicate an agent's private state. We implement this within a hybrid architecture where an offline "Pragmatic Reward Module" computes this reward to augment the agent's learning updates the agent's learning updates, shaping its policy towards legibility without adding latency to its decision-making. This approach provides a formal method for encouraging complex, legible behaviors, bridging the gap between instrumental planning and communicative intent. Our framework provides a formal model for agents that are not only effective but also strategically collaborative.
\end{abstract}
%
\begin{keywords}
Reinforcement Learning, Intrinsic Motivation, Multi-Agent Systems, Computational Pragmatics, Legibility
\end{keywords}
%
\section{Introduction}
\label{sec:intro}

Human collaboration is fundamentally a goal-oriented, strategic activity. To succeed, we don't just act to achieve our own goals; we act in ways that allow our partners to infer our beliefs and intentions, enabling better coordination. The same is true for artificial agents. In many multi-agent reinforcement learning (MARL) environments, like the classic predator-prey game ``Simple Tag,'' an agent's success depends not only on its own actions but also on being understood by its teammates. An agent that simply rushes toward its target may take a path that is ambiguous to its partners, leading to inefficient pursuit. A more sophisticated agent might take a slightly suboptimal path to more clearly ``corral'' the target, making its intention clear and enabling better team coordination.

This property of being easily understood can be called ``legibility.'' However, standard RL agents, optimized solely for an extrinsic task reward, have no incentive to be legible. This leads to our central question: How can we formally encourage an RL agent to choose actions that are not just extrinsically optimal for the task, but also intrinsically optimal for communicating its private state or intentions?

This work proposes a solution in the form of an intrinsic reward, $R_{comm}$, that quantifies the communicative clarity of an agent's actions. To calculate this reward, we model an ``observer'' (e.g., another agent) who uses pragmatic inference to reason about why the agent chose a particular action over the available alternatives. We find the ideal formalism for this observer model in the Rational Speech Act (RSA) framework \cite{goodman2016pragmatic}, which provides a formal account of reasoning about communicative intent.

Our primary contribution is a framework that augments RL with a 'Pragmatic Reward Module.' We build this by first establishing a formal equivalence between the objective functions of Maximum Entropy Reinforcement Learning (Max-Ent RL) and the Rational Speech Act (RSA) framework. This theoretical bridge allows us to derive a principled, intrinsic reward for legible behavior, enabling agents to learn complex, forward-looking strategies that strategically balance task success with communicative clarity.


\section{Related Work}
\label{sec:related}

Our framework, which uses pragmatic inference to generate an intrinsic reward for communicative clarity, builds upon and distinguishes itself from three main areas of research: (1) goal recognition and legibility, which focus on inferring intent from actions; (2) influence-based intrinsic rewards, which incentivize agents to alter their partners' behavior; and (3) emergent communication, which learns explicit messaging protocols.

\subsection{Legibility and Goal Recognition}

The task of inferring an agent's latent state (e.g., its goal or intention) from its observed behavior is the central problem for fields like Inverse Reinforcement Learning (IRL) \cite{ziebart2008maximum} and computational Theory of Mind (ToM) \cite{baker2017rational}. These approaches, like the observer model within our framework, reason backward from an agent's actions to its underlying motivations. Our work is particularly aligned with the concept of "legibility" in robotics, defined as the ability of a robot's motion to make its goal clear to a human observer \cite{dragan2013legibility}.

However, the fundamental difference lies in the purpose and application of this inference. In standard IRL and ToM, the goal is for an external observer to accurately decode an agent's intent. In contrast, our framework places this inferential model inside the learning loop of the acting agent itself. The goal is not for the agent to understand others, but for the agent to generate an internal, intrinsic reward that prospectively shapes its own policy to be more understandable to others. We sidestep the difficult problem of learning an accurate model of another agent's mind by operating under a Centralized Training with Decentralized Execution (CTDE) paradigm. During centralized training, our Pragmatic Reward Module has perfect access to all agents' private states, allowing it to compute a "gold-standard" legibility reward that guides the agent toward communicatively clear policies.

A complementary approach is to use ToM as a source of intrinsic motivation \cite{oguntola2023tom}. Their method provides an intrinsic reward based on an agent's ability to accurately predict the beliefs of other agents. This incentivizes agents not only to learn to model their partners but also to act in ways that make their own beliefs more predictable. This represents a listener-centric approach, rewarding agents for being predictive of others. In contrast, our pragmatic framework is speaker-centric, rewarding agents for being legible and making their own private state clear, regardless of their partners' ability to model it.

\subsection{Influence-Based Intrinsic Rewards}

Methodologically, our pragmatic reward is a form of intrinsic motivation for multi-agent cooperation. This places it in the same category as rewards based on social influence. Notably, Jaques et al. (2019) propose an intrinsic reward that motivates an agent to perform actions that have a "causal influence over other agents' actions" \cite{socialinfluence2019}. This is measured by how much an agent's action changes the probability distribution of its teammates' subsequent actions.

Our approach differs in its fundamental communicative goal. Social influence explicitly rewards an agent for compelling a change in its partners' behavior. Our pragmatic reward, derived from RSA, incentivizes the successful transmission of information about oneself, regardless of how that information is ultimately used. An agent is rewarded for acting in a way that allows an observer to more accurately infer its true private state $(m_t)$. This decouples the speaker's reward from the listener's reaction, allowing the listening agent to incorporate the inferred meaning into its own decision-making or ignore it. This distinction promotes a more flexible form of coordination, where agents reveal their intentions but do not directly force their teammates to comply.

\subsection{Implicit Communication via Action vs. Explicit Communication Channels}

A significant body of work in MARL has focused on agents that learn to communicate through a dedicated, explicit channel \cite{foerster2016learning, sukhbaatar2016learning}. In these models, agents learn to encode and decode messages (typically discrete symbols or continuous vectors) to share information and coordinate. This process results in an "emergent" communication protocol that is optimized alongside the agents' policies for environmental interaction.

Our framework diverges from this paradigm by focusing on implicit communication. There is no separate communication channel or explicit message passing. Instead, communication is embedded within the pragmatics of an agent's task-oriented actions. An agent "speaks" by moving in the environment, and a teammate "listens" by observing that movement. The pragmatic reward incentivizes the agent to choose actions that are not only effective for the task but also serve as clear, unambiguous signals of intent. This approach is particularly valuable in settings where explicit communication is costly, unavailable, or must be combined with physical action, as is common in human-agent teams.


\section{Formalisms for action and interpretation}
\label{sec:background}

To build a model of legible strategic behavior, we require two formalisms: one to model optimal sequential decision-making under uncertainty, and another to model how an observer infers intent from an action. For the former, we turn to Maximum Entropy Reinforcement Learning (Max-Ent RL); for the latter, the Rational Speech Act (RSA) framework.

\subsection{Maximum Entropy Reinforcement Learning (Max-Ent RL)}
\label{ssec:maxentrl}

Reinforcement Learning (RL) provides the mathematical language for sequential decision-making. Maximum Entropy RL (Max-Ent RL) extends the standard RL objective of maximizing rewards with a simple intuition: ``be as versatile as possible while still being successful''. Instead of learning a single, deterministic way to solve a problem, the agent is encouraged to explore multiple near-optimal strategies. Max-Ent RL's key strength is that this entropy bonus leads to more robust and generalizable policies.

In a single-step problem, a Max-Ent RL agent seeks a policy $\pi(a|s)$ that maximizes an objective function comprising both the expected reward and the policy's entropy $H(\pi)$ \cite{haarnoja2018soft}:
\begin{equation}
J(\pi)=\mathbb{E}_{a\sim\pi(\cdot|s)}[R(s,a)]+\alpha_{RL}H(\pi(\cdot|s))
\end{equation}
The optimal policy that maximizes this objective is a softmax distribution over the rewards \cite{levine2018reinforcement}, where the temperature parameter $\alpha_{RL}$ controls the trade-off between exploiting high-reward actions and exploring others:
\begin{equation}
\pi^{*}(a|s)\propto \exp\left(\frac{R(s,a)}{\alpha_{RL}}\right)
\end{equation}

\subsection{The Rational Speech Act (RSA) Framework}
\label{ssec:rsa}

To define our communicative reward, we turn to the Rational Speech Act (RSA) framework \cite{goodman2016pragmatic}, a computational tool for modeling inference about intent. The core intuition is that an observer can infer an agent's intended meaning or private state not just from the literal content of an action, but by reasoning about why the agent chose that specific action over other alternatives. An observer might think, ``If they had intended X, they would have done Y; since they did Z, they must intend something else''. RSA's primary strength is its ability to formally explain complex pragmatic phenomena by modeling this recursive reasoning process.

From an information-theoretic perspective, RSA seeks to find a communicative policy $S(u|m)$ that optimally balances communicative success with action cost. This is formalized by maximizing a gain function, where the resulting model is a softmax distribution over the utility, or ``listener value,'' of each potential communicative act:
\begin{equation}
S(u|m)\propto \exp(\alpha_{RSA}\cdot V_{L}(u,m))
\end{equation}
where $V_{L}(u,m)$ represents the clarity of a communicative act $u$ for conveying meaning $m$ to an observer, penalized by the act's cost.

\section{Rewarding Legibility}
\label{sec:pragmatic_reward}

To encourage legible behavior, our agent optimizes a composite reward function that balances task objectives with communicative clarity. This section defines this reward by establishing a formal bridge between pragmatic inference and reinforcement learning, describes the architecture for its calculation, and presents the resulting strategic Bellman equation that governs the agent's behavior.

\subsection{The Total Reward: Combining Task and Communication}
\label{ssec:total_reward}

We define the agent's decision-making process as being driven by two distinct reward streams at each timestep $t$.

\begin{itemize}
    \item \textbf{Extrinsic Task Reward ($R_{ext}(s_t, a_t)$):} This is the standard reward provided by the environment for task completion. In an environment like ``Simple Tag,'' this might be a positive reward for tagging an opponent or a negative penalty for being tagged by an adversary.
    
    \item \textbf{Intrinsic Communicative Reward ($R_{comm}(s_t, a_t)$):} This is an intrinsic reward that the agent generates internally. It quantifies how clearly action $a_t$ communicates the agent's private state $m_{S_t}$ (henceforth referred to as its 'meaning' or 'intention,' e.g., its intended target) to an imagined observer.
    
    \item \textbf{Total Reward ($R_{total}$):} The agent's policy is optimized to maximize a weighted sum of the extrinsic and intrinsic rewards, where a hyperparameter $\lambda$ balances the importance of task success against communicative clarity.
    \begin{equation}
        R_{total}(s_t, a_t) = R_{ext}(s_t, a_t) + \lambda R_{comm}(s_t, a_t) \label{eq:total_reward}
    \end{equation}
\end{itemize}
This formulation immediately raises the central question: how do we define the communicative reward, $R_{comm}$, in a principled way?

\subsection{Deriving Communicative Reward via Pragmatic Inference}
\label{ssec:reward_derivation}

To define $R_{comm}$, we look to the domain of computational pragmatics. We propose that a principled reward for clarity can be derived by modeling an observer who infers the agent's private state $m_t$ from its action $a_t$. The Rational Speech Act (RSA) framework provides the ideal formalism for this, quantifying the value an observer gains from a communicative act for the purpose of inferring intent. In RSA, this is known as the ``listener value,'' $V_L$, which we select as our candidate for $R_{comm}$.

To justify this choice, we must connect the disparate domains of pragmatic inference and reinforcement learning. We do so by treating the agent's action $a$ as a communicative "utterance" and its private state $m$ as the "meaning" to be conveyed, as shown in Table~\ref{tab:mapping}.

\begin{table}[h]
\centering
\caption{Mapping between Max-Ent RL and RSA components.}
\label{tab:mapping}
\begin{tabular}{|l|l|}
\hline
\textbf{Max-Ent RL Component} & \textbf{RSA Analogue} \\
\hline
State ($s$) & Context \\
Action ($a$) & Utterance ($u$) \\
Private State/Intention ($m$) & Meaning ($m$) \\
Agent Policy ($\pi(a|s,m)$) & Speaker ($S(u|m)$) \\
Intrinsic Reward ($R_{comm}$) & Listener Value ($V_L$) \\
\hline
\end{tabular}
\end{table}

At first glance, the objective of an RL agent and the goal of a pragmatic speaker seem unrelated. However, a closer look at their mathematical foundations reveals a striking formal equivalence. The optimal policy in a single-step Max-Ent RL problem is a softmax distribution over rewards:
\begin{equation}
\pi^{*}(a|s)\propto \exp\left(\frac{R(s,a)}{\alpha_{RL}}\right) \label{eq:maxent_policy}
\end{equation}
Similarly, the RSA speaker model is a softmax over the listener value, governed by an optimality parameter $\alpha_{RSA}$:
\begin{equation}
S(a|m)\propto \exp(\alpha_{RSA}\cdot V_{L}(a,m)) \label{eq:rsa_policy}
\end{equation}
Given that $\alpha_{RSA}$ is analogous to $1/\alpha_{RL}$, the mathematical forms are identical. This formal alignment is the cornerstone of our framework. It provides the theoretical justification to directly import the well-defined listener value ($V_L$) from computational pragmatics and use it as a principled, intrinsic reward ($R_{comm}$) for our RL agent. We thus define our communicative reward as:
\begin{equation}
    R_{comm}(s_t, a_t) \triangleq V_L(a_t, m_{S_t}) = \log L_t^*(m_{S_t}|a_t) - C(a_t) \label{eq:comm_reward_def}
\end{equation}
Here, $L_t^*(m_{S_t}|a_t)$ is the probability a 'literal listener' would assign to the agent's true private state $m_{S_t}$ after observing action $a_t$, quantifying the action's clarity. $C(a_t)$ is an optional cost associated with the action itself (e.g., energy expenditure), which regularizes the policy. In practice, we set $C(a_t)$ to zero because the "cost" of a legible action is already implicitly handled by the task reward. For example, choosing a less direct but clearer path is naturally penalized by a lower $R_{ext}$, making an explicit cost term redundant.

\subsection{The Pragmatic Reward Module: An Architectural Consequence}
\label{ssec:architecture}

Our theoretical bridge allows us to design a practical hybrid architecture that decouples action selection from communicative analysis. The system consists of two core components: the main RL agent, which learns a policy $\pi(a_t|s_t)$, and a \textit{Pragmatic Reward Module}, which calculates the communicative reward to inform the agent's learning process. Crucially, this pragmatic analysis is performed \textit{offline} during the training phase, not in the online action-selection loop, thus adding no latency to the agent's decision-making. After the agent executes an action $a_t$ and experiences a transition $(s_t, a_t, R_{ext}, s_{t+1})$, this module is queried to compute the intrinsic reward, $R_{comm}(s_t, a_t)$, for that specific transition.

To calculate this reward, the module performs the pragmatic inference described in Section \ref{ssec:reward_derivation}. It internally simulates how a rational observer would interpret the chosen action $a_t$ by contrasting it with the set of alternative actions $a' \in A(s_t)$ that were available in that state. The foundation for this internal simulation is the "literal listener" model, $L_0$, which represents the observer's prior belief about the agent's behavior. The implementation of this $L_0$ model is flexible. For instance, $L_0$ could be:
\begin{itemize}
    \item A simple, computationally cheap heuristic, such as a soft-preference for the shortest path in a navigation task.
    \item The agent's own policy, $\pi$, where legibility is measured against the agent's own habits.
    \item A separately trained Theory of Mind (ToM) network \cite{machinetom}, which explicitly models a partner's belief about the agent.
\end{itemize}
This rich, context-aware reward is then used in the strategic Bellman update (Eq.~\ref{eq:full_bellman}) to train the agent's Q-function, shaping the policy to favor actions that are not just extrinsically effective but also intrinsically legible.

\subsection{The Strategic Bellman Equation for Legible Agents}
\label{ssec:bellman}

With the total reward structure defined and the mechanism for calculating $R_{comm}$ established, we can now formulate the full Bellman equation that governs our strategic, legible agent. The action-value ($Q$) and state-value ($V$) functions incorporate both the composite reward and the entropy bonus. The action-value function is the expected future discounted value of taking action $a_t$ in state $s_t$:
\begin{equation}
Q_{t}^{S}(s_{t},a_{t})\triangleq R_{total}(s_{t},a_{t})+\gamma\mathbb{E}_{s_{t+1}}[V_{t+1}^{S}(s_{t+1})]
\end{equation}
The state-value function is the expected value from state $s_t$ when following policy $S$:
\begin{equation}
V_{t}^{S}(s_{t})\triangleq\mathbb{E}_{a_{t}\sim S_{t}(\cdot|s_{t})}[Q_{t}^{S}(s_{t},a_{t})]+\alpha H(S_{t}(\cdot|s_{t}))
\end{equation}
By substituting the optimal softmax policy into the value function definitions, we arrive at the full Bellman Equation for our strategic agent:
\begin{equation}
\label{eq:full_bellman}
\begin{split}
    Q_t(s_t, a_t) ={}& R_{total}(s_t, a_t) + \\
                     & \gamma \mathbb{E}_{s_{t+1}}\left[\alpha \log \sum_{a_{t+1}} \exp\left(\frac{1}{\alpha}Q_{t+1}(s_{t+1}, a_{t+1})\right)\right]
\end{split}
\end{equation}


\section{Analysis and Application}
\label{sec:analysis}

Having formally defined our strategic agent and its underlying computational architecture, we now analyze its properties. First, we illustrate how our framework encourages complex, forward-looking legible behavior in a concrete multi-agent scenario. Second, we demonstrate the generality of our model by showing that it can be specialized to recover prior work on purely communicative agents.

\subsection{Application: Strategic Legibility in a Multi-Agent Environment}
\label{ssec:strategic_behavior}

The primary contribution of our framework is its ability to formally explain and encourage behaviors that strategically balance task success with communicative clarity. We can illustrate this by revisiting the ``Simple Tag'' multi-agent environment, where a team of ``good'' agents must coordinate to tag a faster ``adversary'' agent.

In this scenario, an agent's private state ($m_t$) is the specific adversary it is currently targeting. A standard RL agent, optimizing solely for the extrinsic task reward ($R_{ext}$), might learn to choose the action that moves it on the most direct path toward its target. However, if this path is equidistant between two adversaries, or if it appears to be intercepting a different adversary from a teammate's perspective, the action is ambiguous. This can lead to inefficient team play, where both agents chase the same target or fail to corner one effectively.

Our agent, in contrast, is also motivated by the intrinsic communicative reward ($R_{comm}$). It will learn that a slightly less direct path—one that clearly ``corrals'' a specific adversary relative to its partner's position—is far more legible. This action may have a slightly lower immediate $R_{ext}$ due to the longer path, but it would receive a high $R_{comm}$ because it resolves ambiguity and makes the agent's intention clear to its teammate.

A strategic agent governed by our Bellman equation with a communicative weight $\lambda > 0$ will choose this clarifying action if the expected long-term gain in coordination and capture probability outweighs the immediate cost of the less direct route. This demonstrates that our framework provides a formal mechanism for learning sophisticated, forward-looking strategies that treat communication not as a separate channel, but as an integral part of task-oriented action. Crucially, the same principle applies if a teammate is a human player; an action that is computationally ambiguous to an AI is often deeply confusing to a person, making this intrinsic reward for legibility a key component for smooth human-agent teaming.

\subsection{Generality: Recovering Communicative Agents as a Special Case}
\label{ssec:recover_crsa}

Our general framework for strategic, legible agents can be specialized to recover the behavior of a purely communicative agent, as modeled by frameworks like the Collaborative Rational Speech Act (CRSA) model \cite{estienne2025collaborative}. This is achieved by making two assumptions: the agent has no extrinsic task goal and it is strategically myopic.

This is formalized by setting the extrinsic reward to zero ($R_{ext}=0$) and the communicative reward weight to one ($\lambda=1$). In this scenario, the total reward becomes identical to the communicative reward: $R_{total} = R_{comm}$. Furthermore, if we assume the agent does not plan for the future by setting the discount factor to zero ($\gamma = 0$), the agent's entire objective collapses to maximizing the immediate communicative reward for the current turn, plus its policy entropy. 

This myopic, communication-only objective is equivalent to the per-turn CRSA gain function, $\mathcal{G}_{CRSA}^{\alpha}$. This result demonstrates that purely communicative models are a specific instance of our more general framework, which extends the principles of pragmatic reasoning from single-turn communicative acts to long-horizon, sequential decision-making where instrumental and communicative goals are intertwined.

\section{Methodology}
\label{sec:methodology}

We evaluate our pragmatic reward framework in a series of multi-agent experiments designed to measure the impact of legibility on team coordination and task performance. To ensure reproducibility and fair comparison, all models are implemented in the BenchMARL library and built upon its standard Multi-Agent Soft Actor-Critic (MASAC) algorithm.

\subsection{Environments}
\label{ssec:environments}

We use two standard cooperative tasks from the Multi-Agent Particle Environments (MPE) suite that require significant coordination. In both environments, each agent possesses a private intention, or ``meaning'' ($m_t$), which is included in its local observation.

The first, Simple Tag, is a predator-prey game where a team of 4 agents must capture 2 faster adversaries. An agent's meaning $m_t$ is its intended target, encoded as a one-hot vector. Success requires the team to avoid redundant pursuits and effectively corner the adversaries.

The second, Simple Spread, features 3 agents and 3 landmarks. The team is rewarded based on proximity to landmarks and penalized for collisions. The task requires agents to divide landmarks among themselves, and an agent's meaning $m_t$ is its target landmark.

\subsection{Models and Baselines}
\label{ssec:models}

We implement our models using the Multi-Agent Soft Actor-Critic (MASAC) algorithm, a natural multi-agent extension of the Max-Ent RL framework upon which our theory is built.To isolate the effect of the pragmatic reward, we compare two models. Our proposed model, Pragmatic-MASAC, trains agents to maximize a weighted sum of the extrinsic task reward and our intrinsic pragmatic reward ($R_{total} = R_{ext} + \lambda R_{comm}$). This is compared against a MASAC baseline, where agents are standard learners that optimize solely for the extrinsic task reward, which is equivalent to setting $\lambda=0$. This comparison allows us to directly measure the performance benefits gained from incentivizing legibility.

\subsection{Experimental Conditions}
\label{ssec:conditions}

A core claim of our work is that legible actions allow teammates to infer intent, even without an explicit communication channel. We test this by running all experiments under two conditions that vary the information available to each agent. In the first condition, each agent's observation is augmented with its own private intention ($m_t$) as well as its belief distribution over the intentions of its teammates, providing a direct signal for coordination. In the second, more challenging condition, each agent's observation includes its own private intention but \textit{not} any beliefs about its teammates. Success under this latter condition would provide strong evidence for our framework's utility, as it would demonstrate that agents can learn to infer their teammates' goals purely from their observed actions when incentivized to be legible, removing the need for an explicit belief-sharing mechanism.

\subsection{Evaluation Metrics}
\label{ssec:metrics}

To evaluate our framework, we use a multi-faceted approach. In order to assess team performance, we measure the mean extrinsic reward per episode, averaged over multiple random seeds, comparing the final performance and sample efficiency of Pragmatic-MASAC against a vanilla MASAC baseline. To determine if performance gains are driven by clearer communication, we directly quantify legibility by computing the Average Cross-Entropy between an agent's true intention and an observer's posterior belief. This belief distribution, $P(m_i | a_i, s_t)$, is formed using the listener model from our framework, and a lower cross-entropy value relative to the agent's true one-hot intention vector,$m_i$, indicates that the intent was more accurately and confidently inferred.

\section{Results}
\label{sec:results}

TBD

\section{Discussion}
\label{sec:discussion}

TBD

\subsection{Limitations and Future Work}

While our framework provides a principled method for encouraging legibility, it has several limitations that suggest avenues for future research. A potential confound in our experimental comparison is that the intrinsic reward $R_{comm}$ acts as a form of reward shaping, making the total reward signal less sparse. This may facilitate learning independently of the legibility incentive, complicating direct comparisons of sample efficiency with baseline methods. The framework's utility is also highly dependent on the definitions of an agent's 'meaning' ($m_t$) and the 'literal listener' model ($L_0$). Defining a tractable and useful meaning-space is non-trivial, and we found that different choices for the $L_0$ model can lead to significant variations in performance. One potential solution, using a pre-trained extrinsic-only policy for $L_0$, is computationally expensive and offers no guarantees of being a suitable prior for inference. Furthermore, our approach relies on a Centralized Training with Decentralized Execution (CTDE) paradigm, limiting applicability in fully decentralized settings. Finally, performance is sensitive to key hyperparameters, particularly the reward weight $\lambda$ and the policy temperature $\alpha$, which were not systematically analyzed.

These limitations point to several exciting directions for future work. A primary avenue is to explore more dynamic and complex representations of meaning. For example, meanings could represent an agent's transient intent at a specific turn, beliefs about its surroundings to be implicitly communicated to partners, or be sampled from a large meaning-space to handle real-world complexity. This would expand the agent's communicative capacity beyond the simple intent-signaling explored here. Conversely, our framework could be adapted for adversarial settings by inverting the pragmatic reward to train agents to \textit{obfuscate} their intentions from opponents, while perhaps remaining legible to allies. Further research should also focus on making the framework more autonomous by allowing agents to learn their own belief representations for the $L_0$ model. Lastly, developing methods to make key hyperparameters like $\lambda$ and $\alpha$ learnable or dynamically adjusted could lead to more robust and adaptive agent behavior across different tasks and phases of execution.

% References should be produced using the bibtex program from suitable
% BiBTeX files (here: strings, refs, manuals). The IEEEbib.bst bibliography
% style file from IEEE produces unsorted bibliography list.
% -------------------------------------------------------------------------

\bibliographystyle{IEEEbib}

\begin{thebibliography}{1}

\bibitem{goodman2016pragmatic}
Noah~D. Goodman and Michael~C. Frank,
\newblock ``Pragmatic language interpretation as probabilistic inference,''
\newblock {\em Trends in Cognitive Sciences}, vol. 20, no. 11, pp. 818--829,
  2016.

\bibitem{haarnoja2018soft}
Tuomas Haarnoja, Aurick Zhou, Pieter Abbeel, and Sergey Levine,
\newblock ``Soft actor-critic: Off-policy maximum entropy deep reinforcement
  learning with a stochastic actor,''
\newblock in {\em International Conference on Machine Learning}, 2018, pp.
  1861--1870.

\bibitem{levine2018reinforcement}
Sergey Levine,
\newblock ``Reinforcement learning and control as probabilistic inference:
  Tutorial and review,''
\newblock {\em arXiv preprint arXiv:1805.00909}, 2018.

\bibitem{estienne2025collaborative}
Lautaro Estienne, Gabriel Ben Zenou, Nona Naderi, Jackie Cheung, and Pablo
  Piantanida,
\newblock ``Collaborative rational speech act: Pragmatic reasoning for
  multi-turn dialog,''
\newblock 2025,
\newblock (To appear).

\bibitem{zaslavsky2021rate}
Noga Zaslavsky, Jennifer Hu, and Roger~P. Levy,
\newblock ``A rate-distortion view of human pragmatic reasoning''
\newblock in {\em Proceedings of the Society for Computation in Linguistics
  2021}, 2021, pp. 347--348.

\bibitem{socialinfluence2019}
Jaques, Natasha and Lazaridou, Angeliki and Hughes, Edward and Gulcehre, Caglar and Ortega, Pedro A. and Strouse, DJ and Leibo, Joel Z. and de Freitas, Nando, 
\newblock ``Social influence as intrinsic motivation for multi-agent deep reinforcement learning,``
\newblock 2019

\bibitem{machinetom}
Rabinowitz, N. C., Perbet, F., Song, H. F., Zhang, C., Eslami, S. M. A., and Botvinick, M. 
\newblock ``Machine theory of mind,``
\newblock 2018

\bibitem{ziebart2008maximum}
Ziebart, Brian D and Maas, Andrew L and Bagnell, J Andrew and Dey, Anind K,
\newblock ``Maximum entropy inverse reinforcement learning''
\newblock in {\em Proceedings of the AAAI Conference on Artificial Intelligence, 2008}

\bibitem{baker2017rational}
Baker, Chris L and Jara-Ettinger, Julian and Saxe, Rebecca and Tenenbaum, Joshua B,
\newblock ``Rational quantitative attribution of beliefs, desires, and percepts in human mental state inference'',
\newblock in {\em Advances in Neural Information Processing Systems, 2017}

\bibitem{dragan2013legibility}
Dragan, Anca D and Lee, Kent C and Srinivasa, Siddhartha S,
\newblock ``Legibility and predictability of robot motion'',
\newblock in {\em Interaction Studies vol. 14, 2013}

\bibitem{foerster2016learning}
Foerster, Jakob and Assael, Yannis M and de Freitas, Nando and Whiteson, Shimon,
\newblock ``Learning to communicate with deep multi-agent reinforcement learning'',
\newblock in {\em Advances in Neural Information Processing Systems vol. 29, 2016}

\bibitem{sukhbaatar2016learning}
Sukhbaatar Sainbayar, Szlam Arthur, and Fergus Rob,
\newblock ``Learning Multiagent Communication with Backpropagation'',
\newblock in {\em Advances in Neural Information Processing Systems vol. 29, 2016}

\bibitem{oguntola2023tom}
Oguntola, Ini and Campbell, Joseph and Stepputtis, Simon and Sycara, Katia,
\newblock ``Theory of Mind as Intrinsic Motivation for Multi-Agent Reinforcement Learning'',
\newblock in {\em Proceedings of the First Workshop on Theory of Mind in Communicating Agents vol. 202 , 2023}


\end{thebibliography}


\end{document}